\documentclass{article}

\begin{document}
    \begin{center}
        \Large \textbf{Rabbit Challenge 5: Onboarding und Offboarding} \\ [6pt]
        \normalsize
        Autor: Maximilian Jakob Maag B.Sc. \\
        Version: 1.0
    \end{center}

    \section{Intro}
    Ein neuer Kollege beginnt in Ihrem Team.
    Es ist viel zu tun und die Zeit für die Einarbeitung daher leider sehr kurz.
    Der neue Kollege muss ein Setup erhalten, welches Ihn in die Lage versetzt Infrastrutkur 
    mit Terraform und Ansible zu deployen.

    \section{Aufgaben}
    Ziel dieser Übung ist es Ihr eigenes Setup reproduzierbar zu machen.
    Das soll dabei helfen schneller neue Kollegen onboarden zu können. \\ \\
    Dem neuen Mitarbeiter ist ein Notebook zur Verfügung zu stellen. Es sind alle notwndigen Anwendungen zu installieren und 
    zu konfigurieren.

    \subsection{Aufgabe 1}
    Erstellen Sie eine Dokumentation, welche die Schritte beschreibt, um das Notebook für die Arbeit mit Terraform und Ansible vorzubereiten.
    Diese Dokumentation soll so detailliert sein, dass ein neuer Mitarbeiter ohne Vorkenntnisse in der Lage ist, das Notebook eigenständig einzurichten.
    
    \subsection{Aufgabe 2}
    Implementieren Sie das Setup als bash Script.

    \subsection{Aufgabe 3}
    In Ihrer Abteilung beginnen 10 neue Mitarbeiter. 
    Erstellen Sie ein Ansible Playbook, welches das Setup auf den Notebooks der neuen Mitarbeiter installiert und konfiguriert.

    \subsection{Aufgabe 4}
    Die neuen Kollegen arbeiten nun schon seit einem Jahr in Ihrem Team.
    Das Setup hat sich inzwischen etwas verändert, das neue Setup muss auf alle Notebooks der neuen Kollegen gebracht werden.
    Erstellen Sie ein Ansible Playbook, welches das neue Setup auf den Notebooks der neuen Kollegen installiert und konfiguriert.
    Überlegen Sie, wie Sie fortlaufende Änderungen am Setup in Zukunft einfacher und schneller auf die Notebooks der Kollegen bringen können.

    \subsection{Aufgabe 5}
    Ein Jumphost ist für die Arbeit mit Terraform und Ansible notwendig.
    Diese Änderung wurde nach einem Hackerangriff durch den IT-Sicherheitsbeauftragten angeordnet.
    Deployen Sie einen Jumphost. Mitarbeiter müssen mit einem signierten Public Key über den Jumphost Linux Maschinen in der Cloud erreichen können.

    \subsection{Aufgabe 6}
    Ein Mitarbeiter verlässt das Unternehmen. Alle Zugänge zu Cloud Ressourcen müssen entfernt werden.
    Erstellen Sie ein Ansible Playbook, welches die Zugänge des Mitarbeiters zu Cloud Ressourcen entfernt.
    Die Zugänge der anderen Mitarbeiter müssen dabei unberührt bleiben.

\end{document}